\documentclass{report}
\usepackage{amsmath,amssymb}
\setlength{\parindent}{0mm}
\setlength{\parskip}{1em}
\begin{document}
\begin{center}
\rule{6in}{1pt} \
{\large Bobby Philip \\
{\bf Iterative Methods for Nonlinear Systems Arising in Diblock CoPolymer Systems}}

Computer Science and Mathematics Division \\ Oak Ridge National Laboratory \\ Oak Ridge \\ TN 37831
\\
{\tt bobby.philip@outlook.com}\\
Kumar Rajeev\end{center}

Over the past four decades, extensive scientific research has been
directed towards developing a molecular level understanding of
micro-phase separation in diblock copolymers. Various novel applications
such as the development of advanced materials for thermoplastic
elastomers, nano lithography, fuel cells and supercapacitors have made
these systems a major focus in soft matter research worldwide. Field
theoretic approaches such as the self-consistent field theory (SCFT) have
been used to predict thermodynamic stabilities of various polymer
morphologies such as the lamellae, spheres, cylinders, gyroid etc.

In this talk the SCFT model for a diblock copolymer system is first
described. The resulting set of nonlinear SCFT equations is solved in 3D
using a nonlinear Krylov accelerator solver. Issues encountered in using
potentially more efficient Jacobian Free Newton Krylov (JFNK) solvers and
the performance of less efficient simple mixing algorithms are also
detailed. The outer nonlinear solver requires computing the nonlinear
residual at each iteration. For the SCFT systems this turns out to be
extremely expensive as it requires time integration of a set of diffusion
equations at each iteration. These equations are integrated using an
implicit BDF4 method with multigrid solvers at each time step. After
describing the full solution process some initial results in 3D will be
presented with a description of the challenges to be tackled for future
work.


\end{document}
